% Ngân hàng câu hỏi Toán 10 (KNTT) -- Chương 1, Bài 1, Dạng 01
%% Dạng: Nhận biết mệnh đề, mệnh đề chứa biến; xét tính đúng sai
%% 100 câu: NB 30 -- TH 40 -- VD 20 -- VDC 10 ; mỗi mức: trắc nghiệm : đúng/sai : trả lời ngắn = 3 : 1 : 1
%% Phương án ĐẦU TIÊN trong \cauTN luôn là đáp án đúng (gói trondethi tự đảo).
%% Phần II: \y{nội dung}{d|s}. Phần III: đáp án đặt ở đối số cuối.

\cauTN[id=T10-C1-B1-D01-NB-001,muc=NB,loigiai={Mệnh đề là câu khẳng định có tính đúng hoặc sai xác định. Câu ``$\pi>3{,}2$'' là mệnh đề; các câu còn lại là câu hỏi, câu cảm thán, câu cầu khiến hoặc mệnh đề chứa biến.}]{Câu nào sau đây là mệnh đề?}
  {$\pi>3{,}2$.}
  {Em hãy tính giá trị của biểu thức $2^5$.}{Bạn sinh năm bao nhiêu?}{Đội bóng của trường mình đá hay quá!}

\cauTN[id=T10-C1-B1-D01-NB-002,muc=NB,loigiai={Câu ``Bạn đã làm xong bài tập về nhà chưa?'' là câu hỏi, câu cảm thán hoặc câu cầu khiến nên không phải là mệnh đề. Ba câu còn lại đều là các câu khẳng định có tính đúng sai xác định.}]{Câu nào sau đây \textbf{không} phải là mệnh đề?}
  {Bạn đã làm xong bài tập về nhà chưa?}
  {$\pi>3{,}14$.}{$69$ chia hết cho $7$.}{Số $\pi$ là một số vô tỉ.}

\cauTN[id=T10-C1-B1-D01-NB-003,muc=NB,loigiai={Mệnh đề ``Đường kính là dây cung lớn nhất của đường tròn'' đúng. Kiểm tra trực tiếp, ba mệnh đề còn lại đều sai.}]{Mệnh đề nào sau đây là mệnh đề \textbf{đúng}?}
  {Đường kính là dây cung lớn nhất của đường tròn.}
  {$\pi>3{,}142$.}{Tổng ba góc trong của một tam giác bằng $360^\circ$.}{Hình thang có hai cạnh bên bằng nhau thì là hình thang cân.}

\cauTN[id=T10-C1-B1-D01-NB-004,muc=NB,loigiai={Mệnh đề ``Tổng ba góc trong của một tam giác bằng $360^\circ$'' sai. Ba mệnh đề còn lại đều đúng.}]{Mệnh đề nào sau đây là mệnh đề \textbf{sai}?}
  {Tổng ba góc trong của một tam giác bằng $360^\circ$.}
  {Một năm thường có $365$ ngày.}{Tứ giác có bốn góc vuông là hình chữ nhật.}{Số $0$ là một số tự nhiên.}

\cauTN[id=T10-C1-B1-D01-NB-005,muc=NB,loigiai={Câu ``$n^2+n$ là số chẵn.'' chứa biến, tính đúng sai phụ thuộc vào giá trị của biến nên là mệnh đề chứa biến.}]{Câu nào sau đây là mệnh đề chứa biến?}
  {$n^2+n$ là số chẵn.}
  {Hãy vẽ một tam giác đều cạnh $3$ cm.}{$\sqrt{9}$ là số vô tỉ.}{Hình thang có hai cạnh bên bằng nhau thì là hình thang cân.}

\cauTN[id=T10-C1-B1-D01-NB-006,muc=NB,loigiai={Mệnh đề là câu khẳng định có tính đúng hoặc sai xác định. Câu ``Phương trình $x^2+4x+7=0$ có nghiệm thực'' là mệnh đề; các câu còn lại là câu hỏi, câu cảm thán, câu cầu khiến hoặc mệnh đề chứa biến.}]{Câu nào sau đây là mệnh đề?}
  {Phương trình $x^2+4x+7=0$ có nghiệm thực.}
  {Chúc mừng sinh nhật bạn!}{Đừng nói chuyện riêng trong giờ học!}{Bạn có thích học môn Toán không?}

\cauTN[id=T10-C1-B1-D01-NB-007,muc=NB,loigiai={Câu ``Hôm nay là thứ mấy?'' là câu hỏi, câu cảm thán hoặc câu cầu khiến nên không phải là mệnh đề. Ba câu còn lại đều là các câu khẳng định có tính đúng sai xác định.}]{Câu nào sau đây \textbf{không} phải là mệnh đề?}
  {Hôm nay là thứ mấy?}
  {$16$ là ước của $35$.}{$45$ là số chính phương.}{$\pi>3{,}142$.}

\cauTN[id=T10-C1-B1-D01-NB-008,muc=NB,loigiai={Mệnh đề ``Tam giác đều có ba góc bằng nhau'' đúng. Kiểm tra trực tiếp, ba mệnh đề còn lại đều sai.}]{Mệnh đề nào sau đây là mệnh đề \textbf{đúng}?}
  {Tam giác đều có ba góc bằng nhau.}
  {Tam giác cân luôn có ba cạnh bằng nhau.}{Số $0$ là số nguyên dương.}{$59$ chia hết cho $7$.}

\cauTN[id=T10-C1-B1-D01-NB-009,muc=NB,loigiai={Mệnh đề ``Số $1$ là số nguyên tố'' sai. Ba mệnh đề còn lại đều đúng.}]{Mệnh đề nào sau đây là mệnh đề \textbf{sai}?}
  {Số $1$ là số nguyên tố.}
  {Tổng bốn góc trong của một tứ giác lồi bằng $360^\circ$.}{Phương trình $x^2-3x-5=0$ có nghiệm thực.}{Hình thoi có hai đường chéo vuông góc với nhau.}

\cauTN[id=T10-C1-B1-D01-NB-010,muc=NB,loigiai={Câu ``$x+3=7$.'' chứa biến, tính đúng sai phụ thuộc vào giá trị của biến nên là mệnh đề chứa biến.}]{Câu nào sau đây là mệnh đề chứa biến?}
  {$x+3=7$.}
  {Số $0$ là số nguyên dương.}{$\pi>3{,}14$.}{$\pi$ là một số hữu tỉ.}

\cauTN[id=T10-C1-B1-D01-NB-011,muc=NB,loigiai={Mệnh đề là câu khẳng định có tính đúng hoặc sai xác định. Câu ``Hai tam giác có diện tích bằng nhau thì bằng nhau'' là mệnh đề; các câu còn lại là câu hỏi, câu cảm thán, câu cầu khiến hoặc mệnh đề chứa biến.}]{Câu nào sau đây là mệnh đề?}
  {Hai tam giác có diện tích bằng nhau thì bằng nhau.}
  {Hãy vẽ một tam giác đều cạnh $3$ cm.}{$3x-5<0$.}{Mấy giờ rồi nhỉ?}

\cauTN[id=T10-C1-B1-D01-NB-012,muc=NB,loigiai={Câu ``Em hãy giải phương trình $x^2-1=0$.'' là câu hỏi, câu cảm thán hoặc câu cầu khiến nên không phải là mệnh đề. Ba câu còn lại đều là các câu khẳng định có tính đúng sai xác định.}]{Câu nào sau đây \textbf{không} phải là mệnh đề?}
  {Em hãy giải phương trình $x^2-1=0$.}
  {Mọi số nguyên tố đều là số lẻ.}{Tứ giác có bốn góc vuông là hình chữ nhật.}{$\pi>3{,}142$.}

\cauTN[id=T10-C1-B1-D01-NB-013,muc=NB,loigiai={Mệnh đề ``Phương trình $x^2+2x-2=0$ có nghiệm thực'' đúng. Kiểm tra trực tiếp, ba mệnh đề còn lại đều sai.}]{Mệnh đề nào sau đây là mệnh đề \textbf{đúng}?}
  {Phương trình $x^2+2x-2=0$ có nghiệm thực.}
  {Tam giác cân luôn có ba cạnh bằng nhau.}{Ba số $4$, $5$, $6$ là độ dài ba cạnh của một tam giác vuông.}{$\pi>3{,}15$.}

\cauTN[id=T10-C1-B1-D01-NB-014,muc=NB,loigiai={Mệnh đề ``Hai tam giác có diện tích bằng nhau thì bằng nhau'' sai. Ba mệnh đề còn lại đều đúng.}]{Mệnh đề nào sau đây là mệnh đề \textbf{sai}?}
  {Hai tam giác có diện tích bằng nhau thì bằng nhau.}
  {$\pi>3$.}{Phương trình $x^2+5x+6=0$ có nghiệm thực.}{Hình vuông có bốn cạnh bằng nhau.}

\cauTN[id=T10-C1-B1-D01-NB-015,muc=NB,loigiai={Câu ``$y-2x=1$.'' chứa biến, tính đúng sai phụ thuộc vào giá trị của biến nên là mệnh đề chứa biến.}]{Câu nào sau đây là mệnh đề chứa biến?}
  {$y-2x=1$.}
  {Tứ giác có bốn góc vuông là hình chữ nhật.}{Hôm nay là thứ mấy?}{Ba số $2$, $3$, $4$ là độ dài ba cạnh của một tam giác vuông.}

\cauTN[id=T10-C1-B1-D01-NB-016,muc=NB,loigiai={Mệnh đề là câu khẳng định có tính đúng hoặc sai xác định. Câu ``$25$ là số chính phương'' là mệnh đề; các câu còn lại là câu hỏi, câu cảm thán, câu cầu khiến hoặc mệnh đề chứa biến.}]{Câu nào sau đây là mệnh đề?}
  {$25$ là số chính phương.}
  {Bạn đã làm xong bài tập về nhà chưa?}{Bạn có thích học môn Toán không?}{$n^2+n$ là số chẵn.}

\cauTN[id=T10-C1-B1-D01-NB-017,muc=NB,loigiai={Mệnh đề ``Mọi số nguyên đều là số hữu tỉ'' đúng. Kiểm tra trực tiếp, ba mệnh đề còn lại đều sai.}]{Mệnh đề nào sau đây là mệnh đề \textbf{đúng}?}
  {Mọi số nguyên đều là số hữu tỉ.}
  {Hình chữ nhật luôn có hai đường chéo vuông góc với nhau.}{$18$ là ước của $49$.}{Tam giác cân luôn có ba cạnh bằng nhau.}

\cauTN[id=T10-C1-B1-D01-NB-018,muc=NB,loigiai={Mệnh đề ``$42$ là số chính phương'' sai. Ba mệnh đề còn lại đều đúng.}]{Mệnh đề nào sau đây là mệnh đề \textbf{sai}?}
  {$42$ là số chính phương.}
  {Hình vuông có bốn cạnh bằng nhau.}{Hai đường chéo của hình chữ nhật bằng nhau.}{$\dfrac{3}{2}>\dfrac{2}{6}$.}

\cauDS[id=T10-C1-B1-D01-NB-019,muc=NB,loigiai={Kiểm tra trực tiếp từng mệnh đề. a) Sai. b) Đúng. c) Đúng. d) Sai.}]{Xét tính đúng sai của các mệnh đề sau:}{%
  \y{$14$ là số nguyên tố.}{s}
  \y{Mọi hình vuông đều là hình chữ nhật.}{d}
  \y{$\dfrac{3}{7}>\dfrac{1}{5}$.}{d}
  \y{Mọi số thực đều là số hữu tỉ.}{s}}

\cauDS[id=T10-C1-B1-D01-NB-020,muc=NB,loigiai={Mệnh đề là câu khẳng định đúng hoặc sai; câu hỏi, câu cảm thán, câu cầu khiến và câu chứa biến chưa xác định không phải là mệnh đề. a) Đúng. b) Đúng. c) Sai. d) Sai.}]{Trong giờ học, giáo viên viết lên bảng một số câu. Xét tính đúng sai của các khẳng định sau:}{%
  \y{Câu ``Mọi số thực đều là số hữu tỉ.'' là một mệnh đề.}{d}
  \y{Câu ``Ba số $3$, $4$, $5$ là độ dài ba cạnh của một tam giác vuông.'' là một mệnh đề.}{d}
  \y{Câu ``Bài toán này khó quá!'' là một mệnh đề.}{s}
  \y{Câu ``Phương trình $x^2+1=0$ có nghiệm không?'' là một mệnh đề.}{s}}

\cauDS[id=T10-C1-B1-D01-NB-021,muc=NB,loigiai={Kiểm tra trực tiếp từng mệnh đề. a) Sai. b) Đúng. c) Đúng. d) Đúng.}]{Xét tính đúng sai của các mệnh đề sau:}{%
  \y{Số $0$ là số nguyên dương.}{s}
  \y{Phương trình $x^2+6x+6=0$ có nghiệm thực.}{d}
  \y{Ba số $3$, $4$, $5$ là độ dài ba cạnh của một tam giác vuông.}{d}
  \y{Ba số $5$, $12$, $13$ là độ dài ba cạnh của một tam giác vuông.}{d}}

\cauDS[id=T10-C1-B1-D01-NB-022,muc=NB,loigiai={Mệnh đề là câu khẳng định đúng hoặc sai; câu hỏi, câu cảm thán, câu cầu khiến và câu chứa biến chưa xác định không phải là mệnh đề. a) Đúng. b) Sai. c) Đúng. d) Sai.}]{Trong giờ học, giáo viên viết lên bảng một số câu. Xét tính đúng sai của các khẳng định sau:}{%
  \y{Câu ``Hình vuông có bốn cạnh bằng nhau.'' là một mệnh đề.}{d}
  \y{Câu ``$x+3=7$.'' là một mệnh đề.}{s}
  \y{Câu ``Đường kính là dây cung lớn nhất của đường tròn.'' là một mệnh đề.}{d}
  \y{Câu ``Số $15$ có phải là số nguyên tố không?'' là một mệnh đề.}{s}}

\cauDS[id=T10-C1-B1-D01-NB-023,muc=NB,loigiai={Kiểm tra trực tiếp từng mệnh đề. a) Sai. b) Sai. c) Sai. d) Sai.}]{Xét tính đúng sai của các mệnh đề sau:}{%
  \y{Một tuần có $8$ ngày.}{s}
  \y{Ba số $5$, $6$, $8$ là độ dài ba cạnh của một tam giác vuông.}{s}
  \y{$\pi>3{,}142$.}{s}
  \y{$\dfrac{3}{4}>\dfrac{9}{9}$.}{s}}

\cauDS[id=T10-C1-B1-D01-NB-024,muc=NB,loigiai={Mệnh đề là câu khẳng định đúng hoặc sai; câu hỏi, câu cảm thán, câu cầu khiến và câu chứa biến chưa xác định không phải là mệnh đề. a) Đúng. b) Đúng. c) Sai. d) Sai.}]{Trong giờ học, giáo viên viết lên bảng một số câu. Xét tính đúng sai của các khẳng định sau:}{%
  \y{Câu ``Tam giác cân luôn có ba cạnh bằng nhau.'' là một mệnh đề.}{d}
  \y{Câu ``$98$ là số nguyên tố.'' là một mệnh đề.}{d}
  \y{Câu ``Em hãy tính giá trị của biểu thức $2^5$.'' là một mệnh đề.}{s}
  \y{Câu ``Đừng nói chuyện riêng trong giờ học!'' là một mệnh đề.}{s}}

\cauTLN[id=T10-C1-B1-D01-NB-025,muc=NB,loigiai={Có 5 câu là mệnh đề; các câu còn lại là câu hỏi, câu cảm thán, câu cầu khiến hoặc mệnh đề chứa biến.}]{Cho các câu sau:
\begin{enumerate}[label=(\arabic*),itemsep=0pt]
\item $x$ chia hết cho $5$.
\item $9$ là ước của $116$.
\item Phương trình $x^2-3=0$ có nghiệm thực.
\item $88$ là số nguyên tố.
\item Tứ giác có bốn góc vuông là hình chữ nhật.
\item $|x|=3$.
\item $\dfrac{2}{11}>\dfrac{8}{3}$.
\end{enumerate}
Có bao nhiêu câu là mệnh đề?}{5}

\cauTLN[id=T10-C1-B1-D01-NB-026,muc=NB,loigiai={Các mệnh đề đúng là: (1), (5). Vậy có 2 mệnh đề đúng.}]{Cho các mệnh đề sau:
\begin{enumerate}[label=(\arabic*),itemsep=0pt]
\item Hình thoi có hai đường chéo vuông góc với nhau.
\item $11$ là ước của $87$.
\item $\pi>3{,}2$.
\item $55$ chia hết cho $6$.
\item Đường kính là dây cung lớn nhất của đường tròn.
\end{enumerate}
Có bao nhiêu mệnh đề đúng?}{2}

\cauTLN[id=T10-C1-B1-D01-NB-027,muc=NB,loigiai={Có 4 câu là mệnh đề; các câu còn lại là câu hỏi, câu cảm thán, câu cầu khiến hoặc mệnh đề chứa biến.}]{Cho các câu sau:
\begin{enumerate}[label=(\arabic*),itemsep=0pt]
\item Hôm nay là thứ mấy?
\item Đừng nói chuyện riêng trong giờ học!
\item Mọi hình thoi đều là hình vuông.
\item Đường kính là dây cung lớn nhất của đường tròn.
\item Hình bình hành có hai đường chéo cắt nhau tại trung điểm của mỗi đường.
\item Mọi số nguyên tố đều là số lẻ.
\end{enumerate}
Có bao nhiêu câu là mệnh đề?}{4}

\cauTLN[id=T10-C1-B1-D01-NB-028,muc=NB,loigiai={Các mệnh đề đúng là: (1), (4), (5). Vậy có 3 mệnh đề đúng.}]{Cho các mệnh đề sau:
\begin{enumerate}[label=(\arabic*),itemsep=0pt]
\item $\sqrt{18}$ là số vô tỉ.
\item $\pi$ là một số hữu tỉ.
\item Một tuần có $8$ ngày.
\item Hình bình hành có hai đường chéo cắt nhau tại trung điểm của mỗi đường.
\item Một năm thường có $365$ ngày.
\end{enumerate}
Có bao nhiêu mệnh đề đúng?}{3}

\cauTLN[id=T10-C1-B1-D01-NB-029,muc=NB,loigiai={Có 4 câu là mệnh đề; các câu còn lại là câu hỏi, câu cảm thán, câu cầu khiến hoặc mệnh đề chứa biến.}]{Cho các câu sau:
\begin{enumerate}[label=(\arabic*),itemsep=0pt]
\item Mọi số nguyên đều là số hữu tỉ.
\item Đội bóng của trường mình đá hay quá!
\item $50$ là số nguyên tố.
\item Mọi hình vuông đều là hình chữ nhật.
\item Bạn đã làm xong bài tập về nhà chưa?
\item Hình thoi có hai đường chéo vuông góc với nhau.
\end{enumerate}
Có bao nhiêu câu là mệnh đề?}{4}

\cauTLN[id=T10-C1-B1-D01-NB-030,muc=NB,loigiai={Các mệnh đề đúng là: (1), (3). Vậy có 2 mệnh đề đúng.}]{Cho các mệnh đề sau:
\begin{enumerate}[label=(\arabic*),itemsep=0pt]
\item $\sqrt{2}$ là số vô tỉ.
\item $11$ là ước của $20$.
\item Số $2$ là số nguyên tố chẵn duy nhất.
\item $8$ là ước của $102$.
\item $17$ là ước của $109$.
\end{enumerate}
Có bao nhiêu mệnh đề đúng?}{2}

\cauTN[id=T10-C1-B1-D01-TH-031,muc=TH,loigiai={Thay $x=7$ vào ta được mệnh đề đúng; với các giá trị còn lại ta được mệnh đề sai.}]{Cho mệnh đề chứa biến $P(x)$: ``$3x-16>0$''. Mệnh đề nào sau đây đúng?}
  {$P(7)$.}
  {$P(5)$.}{$P(-3)$.}{$P(1)$.}

\cauTN[id=T10-C1-B1-D01-TH-032,muc=TH,loigiai={Thay $n=0$ vào ta được mệnh đề sai; các giá trị còn lại cho mệnh đề đúng.}]{Cho mệnh đề chứa biến $P(n)$: ``$n^2+2$ chia hết cho $3$''. Mệnh đề nào sau đây \textbf{sai}?}
  {$P(0)$.}
  {$P(-2)$.}{$P(-1)$.}{$P(8)$.}

\cauTN[id=T10-C1-B1-D01-TH-033,muc=TH,loigiai={Kiểm tra từng mệnh đề, có đúng 1 mệnh đề đúng.}]{Cho các mệnh đề:
\begin{enumerate}[label=(\arabic*),itemsep=0pt]
\item Mọi số nguyên tố đều là số lẻ.
\item $\pi>3{,}1$.
\item Ba số $4$, $5$, $6$ là độ dài ba cạnh của một tam giác vuông.
\item $111$ chia hết cho $5$.
\end{enumerate}
Số mệnh đề đúng là}
  {$1$.}
  {$0$.}{$4$.}{$2$.}

\cauTN[id=T10-C1-B1-D01-TH-034,muc=TH,loigiai={Thay $n=7$ vào ta được mệnh đề đúng; với các giá trị còn lại ta được mệnh đề sai.}]{Cho mệnh đề chứa biến $P(n)$: ``$2^n>n^2+3$''. Mệnh đề nào sau đây đúng?}
  {$P(7)$.}
  {$P(3)$.}{$P(-4)$.}{$P(2)$.}

\cauTN[id=T10-C1-B1-D01-TH-035,muc=TH,loigiai={Thay $x=4$ vào ta được mệnh đề sai; các giá trị còn lại cho mệnh đề đúng.}]{Cho mệnh đề chứa biến $P(x)$: ``$x$ là ước nguyên dương của $18$''. Mệnh đề nào sau đây \textbf{sai}?}
  {$P(4)$.}
  {$P(6)$.}{$P(1)$.}{$P(2)$.}

\cauTN[id=T10-C1-B1-D01-TH-036,muc=TH,loigiai={Kiểm tra từng mệnh đề, có đúng 1 mệnh đề đúng.}]{Cho các mệnh đề:
\begin{enumerate}[label=(\arabic*),itemsep=0pt]
\item $29$ là số chính phương.
\item Ba số $5$, $6$, $8$ là độ dài ba cạnh của một tam giác vuông.
\item $\sqrt{16}$ là số vô tỉ.
\item Ba số $3$, $4$, $5$ là độ dài ba cạnh của một tam giác vuông.
\end{enumerate}
Số mệnh đề đúng là}
  {$1$.}
  {$2$.}{$4$.}{$0$.}

\cauTN[id=T10-C1-B1-D01-TH-037,muc=TH,loigiai={Thay $x=1$ vào ta được mệnh đề đúng; với các giá trị còn lại ta được mệnh đề sai.}]{Cho mệnh đề chứa biến $P(x)$: ``$x$ là ước nguyên dương của $12$''. Mệnh đề nào sau đây đúng?}
  {$P(1)$.}
  {$P(0)$.}{$P(7)$.}{$P(5)$.}

\cauTN[id=T10-C1-B1-D01-TH-038,muc=TH,loigiai={Thay $x=-2$ vào ta được mệnh đề sai; các giá trị còn lại cho mệnh đề đúng.}]{Cho mệnh đề chứa biến $P(x)$: ``$5x-10>0$''. Mệnh đề nào sau đây \textbf{sai}?}
  {$P(-2)$.}
  {$P(3)$.}{$P(6)$.}{$P(8)$.}

\cauTN[id=T10-C1-B1-D01-TH-039,muc=TH,loigiai={Kiểm tra từng mệnh đề, có đúng 1 mệnh đề đúng.}]{Cho các mệnh đề:
\begin{enumerate}[label=(\arabic*),itemsep=0pt]
\item $50$ là số chính phương.
\item Ba số $9$, $12$, $15$ là độ dài ba cạnh của một tam giác vuông.
\item $41$ là số chính phương.
\item Mọi số nguyên tố đều là số lẻ.
\end{enumerate}
Số mệnh đề đúng là}
  {$1$.}
  {$2$.}{$3$.}{$0$.}

\cauTN[id=T10-C1-B1-D01-TH-040,muc=TH,loigiai={Thay $x=3$ vào ta được mệnh đề đúng; với các giá trị còn lại ta được mệnh đề sai.}]{Cho mệnh đề chứa biến $P(x)$: ``$x^2-8x+15=0$''. Mệnh đề nào sau đây đúng?}
  {$P(3)$.}
  {$P(1)$.}{$P(2)$.}{$P(4)$.}

\cauTN[id=T10-C1-B1-D01-TH-041,muc=TH,loigiai={Thay $n=0$ vào ta được mệnh đề sai; các giá trị còn lại cho mệnh đề đúng.}]{Cho mệnh đề chứa biến $P(n)$: ``$2^n>n^2+2$''. Mệnh đề nào sau đây \textbf{sai}?}
  {$P(0)$.}
  {$P(7)$.}{$P(8)$.}{$P(6)$.}

\cauTN[id=T10-C1-B1-D01-TH-042,muc=TH,loigiai={Kiểm tra từng mệnh đề, có đúng 3 mệnh đề đúng.}]{Cho các mệnh đề:
\begin{enumerate}[label=(\arabic*),itemsep=0pt]
\item Tứ giác có bốn góc vuông là hình chữ nhật.
\item Ba số $2$, $3$, $4$ là độ dài ba cạnh của một tam giác vuông.
\item Hình vuông có bốn cạnh bằng nhau.
\item Tam giác đều có ba góc bằng nhau.
\end{enumerate}
Số mệnh đề đúng là}
  {$3$.}
  {$4$.}{$1$.}{$0$.}

\cauTN[id=T10-C1-B1-D01-TH-043,muc=TH,loigiai={Thay $n=5$ vào ta được mệnh đề đúng; với các giá trị còn lại ta được mệnh đề sai.}]{Cho mệnh đề chứa biến $P(n)$: ``$2^n>n^2+1$''. Mệnh đề nào sau đây đúng?}
  {$P(5)$.}
  {$P(2)$.}{$P(0)$.}{$P(-4)$.}

\cauTN[id=T10-C1-B1-D01-TH-044,muc=TH,loigiai={Thay $x=-4$ vào ta được mệnh đề sai; các giá trị còn lại cho mệnh đề đúng.}]{Cho mệnh đề chứa biến $P(x)$: ``$7x-15>0$''. Mệnh đề nào sau đây \textbf{sai}?}
  {$P(-4)$.}
  {$P(8)$.}{$P(6)$.}{$P(4)$.}

\cauTN[id=T10-C1-B1-D01-TH-045,muc=TH,loigiai={Kiểm tra từng mệnh đề, có đúng 2 mệnh đề đúng.}]{Cho các mệnh đề:
\begin{enumerate}[label=(\arabic*),itemsep=0pt]
\item $9$ là ước của $45$.
\item $25$ là số nguyên tố.
\item Số $\pi$ là một số vô tỉ.
\item Hai tam giác có diện tích bằng nhau thì bằng nhau.
\end{enumerate}
Số mệnh đề đúng là}
  {$2$.}
  {$3$.}{$1$.}{$4$.}

\cauTN[id=T10-C1-B1-D01-TH-046,muc=TH,loigiai={Thay $x=6$ vào ta được mệnh đề đúng; với các giá trị còn lại ta được mệnh đề sai.}]{Cho mệnh đề chứa biến $P(x)$: ``$x$ là ước nguyên dương của $18$''. Mệnh đề nào sau đây đúng?}
  {$P(6)$.}
  {$P(8)$.}{$P(4)$.}{$P(-2)$.}

\cauTN[id=T10-C1-B1-D01-TH-047,muc=TH,loigiai={Kiểm tra từng mệnh đề, có đúng 3 mệnh đề đúng.}]{Cho các mệnh đề:
\begin{enumerate}[label=(\arabic*),itemsep=0pt]
\item Ba số $6$, $8$, $10$ là độ dài ba cạnh của một tam giác vuông.
\item $14$ là ước của $92$.
\item Ba số $9$, $12$, $15$ là độ dài ba cạnh của một tam giác vuông.
\item Hình bình hành có hai đường chéo cắt nhau tại trung điểm của mỗi đường.
\end{enumerate}
Số mệnh đề đúng là}
  {$3$.}
  {$1$.}{$2$.}{$0$.}

\cauTN[id=T10-C1-B1-D01-TH-048,muc=TH,loigiai={Thay $x=2$ vào ta được mệnh đề đúng; với các giá trị còn lại ta được mệnh đề sai.}]{Cho mệnh đề chứa biến $P(x)$: ``$5x-7>0$''. Mệnh đề nào sau đây đúng?}
  {$P(2)$.}
  {$P(-2)$.}{$P(-3)$.}{$P(0)$.}

\cauTN[id=T10-C1-B1-D01-TH-049,muc=TH,loigiai={Thay $x=0$ vào ta được mệnh đề sai; các giá trị còn lại cho mệnh đề đúng.}]{Cho mệnh đề chứa biến $P(x)$: ``$3x-13>0$''. Mệnh đề nào sau đây \textbf{sai}?}
  {$P(0)$.}
  {$P(7)$.}{$P(6)$.}{$P(8)$.}

\cauTN[id=T10-C1-B1-D01-TH-050,muc=TH,loigiai={Kiểm tra từng mệnh đề, có đúng 2 mệnh đề đúng.}]{Cho các mệnh đề:
\begin{enumerate}[label=(\arabic*),itemsep=0pt]
\item $90$ chia hết cho $3$.
\item Số $2$ là số nguyên tố chẵn duy nhất.
\item $55$ là số chính phương.
\item $\sqrt{25}$ là số vô tỉ.
\end{enumerate}
Số mệnh đề đúng là}
  {$2$.}
  {$3$.}{$0$.}{$1$.}

\cauTN[id=T10-C1-B1-D01-TH-051,muc=TH,loigiai={Thay $x=-3$ vào ta được mệnh đề sai; các giá trị còn lại cho mệnh đề đúng.}]{Cho mệnh đề chứa biến $P(x)$: ``$x$ là ước nguyên dương của $20$''. Mệnh đề nào sau đây \textbf{sai}?}
  {$P(-3)$.}
  {$P(5)$.}{$P(2)$.}{$P(4)$.}

\cauTN[id=T10-C1-B1-D01-TH-052,muc=TH,loigiai={Kiểm tra từng mệnh đề, có đúng 2 mệnh đề đúng.}]{Cho các mệnh đề:
\begin{enumerate}[label=(\arabic*),itemsep=0pt]
\item $\sqrt{18}$ là số vô tỉ.
\item Tổng bốn góc trong của một tứ giác lồi bằng $360^\circ$.
\item $91$ là số chính phương.
\item $14$ là ước của $30$.
\end{enumerate}
Số mệnh đề đúng là}
  {$2$.}
  {$0$.}{$4$.}{$3$.}

\cauTN[id=T10-C1-B1-D01-TH-053,muc=TH,loigiai={Thay $x=2$ vào ta được mệnh đề đúng; với các giá trị còn lại ta được mệnh đề sai.}]{Cho mệnh đề chứa biến $P(x)$: ``$x^2-4=0$''. Mệnh đề nào sau đây đúng?}
  {$P(2)$.}
  {$P(0)$.}{$P(7)$.}{$P(3)$.}

\cauTN[id=T10-C1-B1-D01-TH-054,muc=TH,loigiai={Thay $x=0$ vào ta được mệnh đề sai; các giá trị còn lại cho mệnh đề đúng.}]{Cho mệnh đề chứa biến $P(x)$: ``$x$ là ước nguyên dương của $12$''. Mệnh đề nào sau đây \textbf{sai}?}
  {$P(0)$.}
  {$P(6)$.}{$P(2)$.}{$P(3)$.}

\cauDS[id=T10-C1-B1-D01-TH-055,muc=TH,loigiai={Thay trực tiếp các giá trị của $n$ để kiểm tra. a) Đúng. b) Đúng. c) Sai. d) Đúng.}]{Cho mệnh đề chứa biến $P(n)$: ``$2^n>n^2+2$'' với $n$ là số nguyên.}{%
  \y{$P(6)$ là mệnh đề đúng.}{d}
  \y{$P(5)$ là mệnh đề đúng.}{d}
  \y{$P(-3)$ là mệnh đề đúng.}{s}
  \y{Tồn tại số tự nhiên $n$ để $P(n)$ là mệnh đề đúng.}{d}}

\cauDS[id=T10-C1-B1-D01-TH-056,muc=TH,loigiai={Thay trực tiếp các giá trị của $n$ để kiểm tra. a) Đúng. b) Sai. c) Sai. d) Đúng.}]{Cho mệnh đề chứa biến $P(n)$: ``$2^n>n^2+2$'' với $n$ là số nguyên.}{%
  \y{$P(6)$ là mệnh đề đúng.}{d}
  \y{$P(-1)$ là mệnh đề đúng.}{s}
  \y{$P(-2)$ là mệnh đề đúng.}{s}
  \y{Tồn tại số tự nhiên $n$ để $P(n)$ là mệnh đề đúng.}{d}}

\cauDS[id=T10-C1-B1-D01-TH-057,muc=TH,loigiai={Thay trực tiếp các giá trị của $n$ để kiểm tra. a) Đúng. b) Sai. c) Đúng. d) Đúng.}]{Cho mệnh đề chứa biến $P(n)$: ``$2^n>n^2+2$'' với $n$ là số nguyên.}{%
  \y{$P(5)$ là mệnh đề đúng.}{d}
  \y{$P(1)$ là mệnh đề đúng.}{s}
  \y{$P(8)$ là mệnh đề đúng.}{d}
  \y{Tồn tại số tự nhiên $n$ để $P(n)$ là mệnh đề đúng.}{d}}

\cauDS[id=T10-C1-B1-D01-TH-058,muc=TH,loigiai={Thay trực tiếp các giá trị của $x$ để kiểm tra. a) Sai. b) Sai. c) Sai. d) Đúng.}]{Cho mệnh đề chứa biến $P(x)$: ``$x$ là ước nguyên dương của $36$'' với $x$ là số nguyên.}{%
  \y{$P(0)$ là mệnh đề đúng.}{s}
  \y{$P(-1)$ là mệnh đề đúng.}{s}
  \y{$P(-2)$ là mệnh đề đúng.}{s}
  \y{Tồn tại số tự nhiên $x$ để $P(x)$ là mệnh đề đúng.}{d}}

\cauDS[id=T10-C1-B1-D01-TH-059,muc=TH,loigiai={Thay trực tiếp các giá trị của $x$ để kiểm tra. a) Đúng. b) Sai. c) Sai. d) Đúng.}]{Cho mệnh đề chứa biến $P(x)$: ``$x$ là ước nguyên dương của $18$'' với $x$ là số nguyên.}{%
  \y{$P(6)$ là mệnh đề đúng.}{d}
  \y{$P(4)$ là mệnh đề đúng.}{s}
  \y{$P(8)$ là mệnh đề đúng.}{s}
  \y{Tồn tại số tự nhiên $x$ để $P(x)$ là mệnh đề đúng.}{d}}

\cauDS[id=T10-C1-B1-D01-TH-060,muc=TH,loigiai={Thay trực tiếp các giá trị của $n$ để kiểm tra. a) Sai. b) Đúng. c) Sai. d) Đúng.}]{Cho mệnh đề chứa biến $P(n)$: ``$2^n>n^2+2$'' với $n$ là số nguyên.}{%
  \y{$P(4)$ là mệnh đề đúng.}{s}
  \y{$P(8)$ là mệnh đề đúng.}{d}
  \y{$P(2)$ là mệnh đề đúng.}{s}
  \y{Tồn tại số tự nhiên $n$ để $P(n)$ là mệnh đề đúng.}{d}}

\cauDS[id=T10-C1-B1-D01-TH-061,muc=TH,loigiai={Thay trực tiếp các giá trị của $n$ để kiểm tra. a) Sai. b) Sai. c) Sai. d) Sai.}]{Cho mệnh đề chứa biến $P(n)$: ``$n^2+2$ chia hết cho $5$'' với $n$ là số nguyên.}{%
  \y{$P(7)$ là mệnh đề đúng.}{s}
  \y{$P(1)$ là mệnh đề đúng.}{s}
  \y{$P(2)$ là mệnh đề đúng.}{s}
  \y{Tồn tại số tự nhiên $n$ để $P(n)$ là mệnh đề đúng.}{s}}

\cauDS[id=T10-C1-B1-D01-TH-062,muc=TH,loigiai={Thay trực tiếp các giá trị của $n$ để kiểm tra. a) Đúng. b) Đúng. c) Sai. d) Đúng.}]{Cho mệnh đề chứa biến $P(n)$: ``$n^2+5$ chia hết cho $6$'' với $n$ là số nguyên.}{%
  \y{$P(1)$ là mệnh đề đúng.}{d}
  \y{$P(5)$ là mệnh đề đúng.}{d}
  \y{$P(6)$ là mệnh đề đúng.}{s}
  \y{Tồn tại số tự nhiên $n$ để $P(n)$ là mệnh đề đúng.}{d}}

\cauTLN[id=T10-C1-B1-D01-TH-063,muc=TH,loigiai={Kiểm tra lần lượt các số nguyên từ $-2$ đến $9$, có 7 giá trị thoả mãn.}]{Cho mệnh đề chứa biến $P(x)$: ``$8x-22>0$''. Có bao nhiêu số nguyên $x$ thuộc đoạn $[-2;9]$ để $P(x)$ là mệnh đề đúng?}{7}

\cauTLN[id=T10-C1-B1-D01-TH-064,muc=TH,loigiai={Kiểm tra lần lượt các số nguyên từ $-2$ đến $7$, có 5 giá trị thoả mãn.}]{Cho mệnh đề chứa biến $P(x)$: ``$x$ là ước nguyên dương của $12$''. Có bao nhiêu số nguyên $x$ thuộc đoạn $[-2;7]$ để $P(x)$ là mệnh đề đúng?}{5}

\cauTLN[id=T10-C1-B1-D01-TH-065,muc=TH,loigiai={Kiểm tra lần lượt các số nguyên từ $-3$ đến $12$, có 8 giá trị thoả mãn.}]{Cho mệnh đề chứa biến $P(n)$: ``$2^n>n^2+4$''. Có bao nhiêu số nguyên $n$ thuộc đoạn $[-3;12]$ để $P(n)$ là mệnh đề đúng?}{8}

\cauTLN[id=T10-C1-B1-D01-TH-066,muc=TH,loigiai={Kiểm tra lần lượt các số nguyên từ $-3$ đến $9$, có 8 giá trị thoả mãn.}]{Cho mệnh đề chứa biến $P(x)$: ``$4x-5>0$''. Có bao nhiêu số nguyên $x$ thuộc đoạn $[-3;9]$ để $P(x)$ là mệnh đề đúng?}{8}

\cauTLN[id=T10-C1-B1-D01-TH-067,muc=TH,loigiai={Kiểm tra lần lượt các số nguyên từ $-4$ đến $9$, có 2 giá trị thoả mãn.}]{Cho mệnh đề chứa biến $P(x)$: ``$x^2-4x=0$''. Có bao nhiêu số nguyên $x$ thuộc đoạn $[-4;9]$ để $P(x)$ là mệnh đề đúng?}{2}

\cauTLN[id=T10-C1-B1-D01-TH-068,muc=TH,loigiai={Kiểm tra lần lượt các số nguyên từ $0$ đến $15$, có 7 giá trị thoả mãn.}]{Cho mệnh đề chứa biến $P(x)$: ``$x$ là ước nguyên dương của $30$''. Có bao nhiêu số nguyên $x$ thuộc đoạn $[0;15]$ để $P(x)$ là mệnh đề đúng?}{7}

\cauTLN[id=T10-C1-B1-D01-TH-069,muc=TH,loigiai={Kiểm tra lần lượt các số nguyên từ $-1$ đến $15$, có 11 giá trị thoả mãn.}]{Cho mệnh đề chứa biến $P(n)$: ``$2^n>n^2+2$''. Có bao nhiêu số nguyên $n$ thuộc đoạn $[-1;15]$ để $P(n)$ là mệnh đề đúng?}{11}

\cauTLN[id=T10-C1-B1-D01-TH-070,muc=TH,loigiai={Kiểm tra lần lượt các số nguyên từ $-3$ đến $13$, có 2 giá trị thoả mãn.}]{Cho mệnh đề chứa biến $P(x)$: ``$x^2-2x-3=0$''. Có bao nhiêu số nguyên $x$ thuộc đoạn $[-3;13]$ để $P(x)$ là mệnh đề đúng?}{2}

\cauTN[id=T10-C1-B1-D01-VD-071,muc=VD,loigiai={Các câu là mệnh đề: ``Số $2$ là số nguyên tố chẵn duy nhất.'' (đúng); ``$13$ là ước của $61$.'' (sai); ``Tam giác đều có ba góc bằng nhau.'' (đúng). Các câu còn lại không phải mệnh đề. Vậy có 2 mệnh đề đúng.}]{Trong buổi sinh hoạt câu lạc bộ Toán học của trường THPT Nguyễn Huệ, bạn Hương được giao nhiệm vụ viết lên bảng năm câu để cả nhóm cùng thảo luận về khái niệm mệnh đề. Các câu bạn Hương viết như sau:
\begin{enumerate}[label=(\arabic*),itemsep=0pt]
\item Tam giác đều có ba góc bằng nhau.
\item Số $2$ là số nguyên tố chẵn duy nhất.
\item Đội bóng của trường mình đá hay quá!
\item Đừng nói chuyện riêng trong giờ học!
\item $13$ là ước của $61$.
\end{enumerate}
Giáo viên hướng dẫn yêu cầu cả nhóm đếm xem trong năm câu trên có bao nhiêu câu là mệnh đề \textbf{đúng}. Kết quả đúng là}
  {$2$.}
  {$0$.}{$4$.}{$5$.}

\cauTN[id=T10-C1-B1-D01-VD-072,muc=VD,loigiai={Số các bội của $4$ không vượt quá $77$ là $19$; trong đó số chia hết cho cả $4$ và $9$ (tức chia hết cho $36$) là $2$. Vậy có $19-2=17$ số.}]{Xét mệnh đề chứa biến $P(n)$: ``$n$ chia hết cho $4$ nhưng không chia hết cho $9$'' với $n$ là số nguyên dương. Có bao nhiêu số nguyên dương $n\le 77$ để $P(n)$ là mệnh đề đúng?}
  {$17$.}
  {$18$.}{$16$.}{$19$.}

\cauTN[id=T10-C1-B1-D01-VD-073,muc=VD,loigiai={Các câu là mệnh đề: ``$\dfrac{8}{5}>\dfrac{8}{9}$.'' (đúng); ``$166$ chia hết cho $4$.'' (sai); ``Một tuần có $8$ ngày.'' (sai). Các câu còn lại không phải mệnh đề. Vậy có 1 mệnh đề đúng.}]{Trong buổi sinh hoạt câu lạc bộ Toán học của trường THPT Lê Lợi, bạn Nam được giao nhiệm vụ viết lên bảng năm câu để cả nhóm cùng thảo luận về khái niệm mệnh đề. Các câu bạn Nam viết như sau:
\begin{enumerate}[label=(\arabic*),itemsep=0pt]
\item Hãy đóng cửa sổ lại!
\item Đội bóng của trường mình đá hay quá!
\item $166$ chia hết cho $4$.
\item $\dfrac{8}{5}>\dfrac{8}{9}$.
\item Một tuần có $8$ ngày.
\end{enumerate}
Giáo viên hướng dẫn yêu cầu cả nhóm đếm xem trong năm câu trên có bao nhiêu câu là mệnh đề \textbf{đúng}. Kết quả đúng là}
  {$1$.}
  {$3$.}{$0$.}{$5$.}

\cauTN[id=T10-C1-B1-D01-VD-074,muc=VD,loigiai={Số các bội của $2$ không vượt quá $76$ là $38$; trong đó số chia hết cho cả $2$ và $3$ (tức chia hết cho $6$) là $12$. Vậy có $38-12=26$ số.}]{Xét mệnh đề chứa biến $P(n)$: ``$n$ chia hết cho $2$ nhưng không chia hết cho $3$'' với $n$ là số nguyên dương. Có bao nhiêu số nguyên dương $n\le 76$ để $P(n)$ là mệnh đề đúng?}
  {$26$.}
  {$27$.}{$25$.}{$38$.}

\cauTN[id=T10-C1-B1-D01-VD-075,muc=VD,loigiai={Các câu là mệnh đề: ``$\dfrac{7}{3}>\dfrac{4}{3}$.'' (đúng); ``$82$ chia hết cho $8$.'' (sai); ``Tổng ba góc trong của một tam giác bằng $360^\circ$.'' (sai). Các câu còn lại không phải mệnh đề. Vậy có 1 mệnh đề đúng.}]{Trong buổi sinh hoạt câu lạc bộ Toán học của trường THPT Trần Hưng Đạo, bạn Thảo được giao nhiệm vụ viết lên bảng năm câu để cả nhóm cùng thảo luận về khái niệm mệnh đề. Các câu bạn Thảo viết như sau:
\begin{enumerate}[label=(\arabic*),itemsep=0pt]
\item Hãy vẽ một tam giác đều cạnh $3$ cm.
\item Tổng ba góc trong của một tam giác bằng $360^\circ$.
\item $x^2>4$.
\item $82$ chia hết cho $8$.
\item $\dfrac{7}{3}>\dfrac{4}{3}$.
\end{enumerate}
Giáo viên hướng dẫn yêu cầu cả nhóm đếm xem trong năm câu trên có bao nhiêu câu là mệnh đề \textbf{đúng}. Kết quả đúng là}
  {$1$.}
  {$5$.}{$2$.}{$4$.}

\cauTN[id=T10-C1-B1-D01-VD-076,muc=VD,loigiai={Số các bội của $2$ không vượt quá $70$ là $35$; trong đó số chia hết cho cả $2$ và $15$ (tức chia hết cho $30$) là $2$. Vậy có $35-2=33$ số.}]{Xét mệnh đề chứa biến $P(n)$: ``$n$ chia hết cho $2$ nhưng không chia hết cho $15$'' với $n$ là số nguyên dương. Có bao nhiêu số nguyên dương $n\le 70$ để $P(n)$ là mệnh đề đúng?}
  {$33$.}
  {$34$.}{$32$.}{$35$.}

\cauTN[id=T10-C1-B1-D01-VD-077,muc=VD,loigiai={Các câu là mệnh đề: ``Một tuần có $8$ ngày.'' (sai); ``Số $0$ là số nguyên dương.'' (sai). Các câu còn lại không phải mệnh đề. Vậy có 0 mệnh đề đúng.}]{Trong buổi sinh hoạt câu lạc bộ Toán học của trường THPT Chu Văn An, bạn Uyên được giao nhiệm vụ viết lên bảng năm câu để cả nhóm cùng thảo luận về khái niệm mệnh đề. Các câu bạn Uyên viết như sau:
\begin{enumerate}[label=(\arabic*),itemsep=0pt]
\item $3x-5<0$.
\item Lớp mình có bao nhiêu bạn nữ?
\item Một tuần có $8$ ngày.
\item Số $0$ là số nguyên dương.
\item Đường từ nhà bạn đến trường dài bao nhiêu kilômét?
\end{enumerate}
Giáo viên hướng dẫn yêu cầu cả nhóm đếm xem trong năm câu trên có bao nhiêu câu là mệnh đề \textbf{đúng}. Kết quả đúng là}
  {$0$.}
  {$1$.}{$4$.}{$5$.}

\cauTN[id=T10-C1-B1-D01-VD-078,muc=VD,loigiai={Số các bội của $6$ không vượt quá $72$ là $12$; trong đó số chia hết cho cả $6$ và $15$ (tức chia hết cho $30$) là $2$. Vậy có $12-2=10$ số.}]{Xét mệnh đề chứa biến $P(n)$: ``$n$ chia hết cho $6$ nhưng không chia hết cho $15$'' với $n$ là số nguyên dương. Có bao nhiêu số nguyên dương $n\le 72$ để $P(n)$ là mệnh đề đúng?}
  {$10$.}
  {$11$.}{$9$.}{$12$.}

\cauTN[id=T10-C1-B1-D01-VD-079,muc=VD,loigiai={Các câu là mệnh đề: ``Số $0$ là số nguyên dương.'' (sai); ``$105$ chia hết cho $10$.'' (sai); ``$77$ là số chính phương.'' (sai). Các câu còn lại không phải mệnh đề. Vậy có 0 mệnh đề đúng.}]{Trong buổi sinh hoạt câu lạc bộ Toán học của trường THPT Nguyễn Huệ, bạn Vy được giao nhiệm vụ viết lên bảng năm câu để cả nhóm cùng thảo luận về khái niệm mệnh đề. Các câu bạn Vy viết như sau:
\begin{enumerate}[label=(\arabic*),itemsep=0pt]
\item $77$ là số chính phương.
\item $105$ chia hết cho $10$.
\item Số $0$ là số nguyên dương.
\item $|x|=3$.
\item Hãy vẽ một tam giác đều cạnh $3$ cm.
\end{enumerate}
Giáo viên hướng dẫn yêu cầu cả nhóm đếm xem trong năm câu trên có bao nhiêu câu là mệnh đề \textbf{đúng}. Kết quả đúng là}
  {$0$.}
  {$1$.}{$4$.}{$5$.}

\cauTN[id=T10-C1-B1-D01-VD-080,muc=VD,loigiai={Số các bội của $3$ không vượt quá $50$ là $16$; trong đó số chia hết cho cả $3$ và $4$ (tức chia hết cho $12$) là $4$. Vậy có $16-4=12$ số.}]{Xét mệnh đề chứa biến $P(n)$: ``$n$ chia hết cho $3$ nhưng không chia hết cho $4$'' với $n$ là số nguyên dương. Có bao nhiêu số nguyên dương $n\le 50$ để $P(n)$ là mệnh đề đúng?}
  {$12$.}
  {$13$.}{$11$.}{$16$.}

\cauTN[id=T10-C1-B1-D01-VD-081,muc=VD,loigiai={Các câu là mệnh đề: ``Phương trình $x^2-3x+4=0$ có nghiệm thực.'' (sai); ``$103$ chia hết cho $8$.'' (sai); ``$46$ là số nguyên tố.'' (sai). Các câu còn lại không phải mệnh đề. Vậy có 0 mệnh đề đúng.}]{Trong buổi sinh hoạt câu lạc bộ Toán học của trường THPT Chu Văn An, bạn Cường được giao nhiệm vụ viết lên bảng năm câu để cả nhóm cùng thảo luận về khái niệm mệnh đề. Các câu bạn Cường viết như sau:
\begin{enumerate}[label=(\arabic*),itemsep=0pt]
\item $\sqrt{x}=2$.
\item Bài toán này khó quá!
\item $103$ chia hết cho $8$.
\item Phương trình $x^2-3x+4=0$ có nghiệm thực.
\item $46$ là số nguyên tố.
\end{enumerate}
Giáo viên hướng dẫn yêu cầu cả nhóm đếm xem trong năm câu trên có bao nhiêu câu là mệnh đề \textbf{đúng}. Kết quả đúng là}
  {$0$.}
  {$4$.}{$2$.}{$1$.}

\cauTN[id=T10-C1-B1-D01-VD-082,muc=VD,loigiai={Số các bội của $4$ không vượt quá $64$ là $16$; trong đó số chia hết cho cả $4$ và $10$ (tức chia hết cho $20$) là $3$. Vậy có $16-3=13$ số.}]{Xét mệnh đề chứa biến $P(n)$: ``$n$ chia hết cho $4$ nhưng không chia hết cho $10$'' với $n$ là số nguyên dương. Có bao nhiêu số nguyên dương $n\le 64$ để $P(n)$ là mệnh đề đúng?}
  {$13$.}
  {$14$.}{$12$.}{$16$.}

\cauDS[id=T10-C1-B1-D01-VD-083,muc=VD,loigiai={Ta có $x^2-2x=(x)(x-2)\le0\Leftrightarrow 0\le x\le 2$. Các số nguyên thoả mãn là $3$ số, có tổng bằng $3$. a) Đúng. b) Đúng. c) Sai. d) Sai.}]{Cho mệnh đề chứa biến $P(x)$: ``$x^2-2x\le 0$'' với $x$ là số nguyên.}{%
  \y{$P(2)$ là mệnh đề đúng.}{d}
  \y{$P(3)$ là mệnh đề sai.}{d}
  \y{Có đúng $4$ số nguyên $x$ để $P(x)$ là mệnh đề đúng.}{s}
  \y{Tổng các số nguyên $x$ để $P(x)$ đúng bằng $3$.}{s}}

\cauDS[id=T10-C1-B1-D01-VD-084,muc=VD,loigiai={Ta có $x^2+x-12=(x+4)(x-3)\le0\Leftrightarrow -4\le x\le 3$. Các số nguyên thoả mãn là $8$ số, có tổng bằng $-4$. a) Đúng. b) Đúng. c) Đúng. d) Sai.}]{Cho mệnh đề chứa biến $P(x)$: ``$x^2+x-12\le 0$'' với $x$ là số nguyên.}{%
  \y{$P(1)$ là mệnh đề đúng.}{d}
  \y{$P(4)$ là mệnh đề sai.}{d}
  \y{Có đúng $8$ số nguyên $x$ để $P(x)$ là mệnh đề đúng.}{d}
  \y{Tổng các số nguyên $x$ để $P(x)$ đúng bằng $-4$.}{s}}

\cauDS[id=T10-C1-B1-D01-VD-085,muc=VD,loigiai={Ta có $x^2-x-12=(x+3)(x-4)\le0\Leftrightarrow -3\le x\le 4$. Các số nguyên thoả mãn là $8$ số, có tổng bằng $4$. a) Sai. b) Đúng. c) Đúng. d) Sai.}]{Cho mệnh đề chứa biến $P(x)$: ``$x^2-x-12\le 0$'' với $x$ là số nguyên.}{%
  \y{$P(-5)$ là mệnh đề đúng.}{s}
  \y{$P(5)$ là mệnh đề sai.}{d}
  \y{Có đúng $8$ số nguyên $x$ để $P(x)$ là mệnh đề đúng.}{d}
  \y{Tổng các số nguyên $x$ để $P(x)$ đúng bằng $5$.}{s}}

\cauDS[id=T10-C1-B1-D01-VD-086,muc=VD,loigiai={Ta có $x^2-1=(x+1)(x-1)\le0\Leftrightarrow -1\le x\le 1$. Các số nguyên thoả mãn là $3$ số, có tổng bằng $0$. a) Đúng. b) Đúng. c) Đúng. d) Sai.}]{Cho mệnh đề chứa biến $P(x)$: ``$x^2-1\le 0$'' với $x$ là số nguyên.}{%
  \y{$P(-1)$ là mệnh đề đúng.}{d}
  \y{$P(2)$ là mệnh đề sai.}{d}
  \y{Có đúng $3$ số nguyên $x$ để $P(x)$ là mệnh đề đúng.}{d}
  \y{Tổng các số nguyên $x$ để $P(x)$ đúng bằng $0$.}{s}}

\cauTLN[id=T10-C1-B1-D01-VD-087,muc=VD,loigiai={Lần lượt thử $n=0,1,2,\ldots$ ta thấy $n=2$ cho $n^2+n+3=9$ không phải số nguyên tố, các giá trị nhỏ hơn đều cho số nguyên tố.}]{Bạn Dũng cho rằng: ``Với mọi số tự nhiên $n$, số $n^2+n+3$ luôn là số nguyên tố''. Để chứng tỏ khẳng định của bạn ấy là sai, cần chỉ ra một phản ví dụ. Tìm số tự nhiên $n$ nhỏ nhất làm phản ví dụ.}{2}

\cauTLN[id=T10-C1-B1-D01-VD-088,muc=VD,loigiai={Đếm các $n$ để $5n+19$ chia hết cho $3$, các $n$ để chia hết cho $5$, trừ đi các $n$ để chia hết cho $15$. Kết quả là $11$.}]{Xét mệnh đề chứa biến $P(n)$: ``$5n+19$ chia hết cho $3$ hoặc chia hết cho $5$''. Có bao nhiêu số nguyên dương $n\le 32$ để $P(n)$ đúng?}{11}

\cauTLN[id=T10-C1-B1-D01-VD-089,muc=VD,loigiai={Đếm các $n$ để $3n+4$ chia hết cho $3$, các $n$ để chia hết cho $5$, trừ đi các $n$ để chia hết cho $15$. Kết quả là $5$.}]{Xét mệnh đề chứa biến $P(n)$: ``$3n+4$ chia hết cho $3$ hoặc chia hết cho $5$''. Có bao nhiêu số nguyên dương $n\le 22$ để $P(n)$ đúng?}{5}

\cauTLN[id=T10-C1-B1-D01-VD-090,muc=VD,loigiai={Lần lượt thử $n=0,1,2,\ldots$ ta thấy $n=4$ cho $n^2+n+5=25$ không phải số nguyên tố, các giá trị nhỏ hơn đều cho số nguyên tố.}]{Bạn Quân cho rằng: ``Với mọi số tự nhiên $n$, số $n^2+n+5$ luôn là số nguyên tố''. Để chứng tỏ khẳng định của bạn ấy là sai, cần chỉ ra một phản ví dụ. Tìm số tự nhiên $n$ nhỏ nhất làm phản ví dụ.}{4}

\cauTN[id=T10-C1-B1-D01-VDC-091,muc=VDC,loigiai={Thử lần lượt: $n=9$ thì $2^{9}=512\le 733$; $n=10$ thì $2^{10}=1024>1004$. Vậy $n=10$.}]{Cho mệnh đề chứa biến $P(n)$: ``$2^n>n^3+4$'' với $n$ là số nguyên dương. Số nguyên dương $n$ nhỏ nhất để $P(n)$ là mệnh đề đúng là}
  {$10$.}
  {$9$.}{$11$.}{$12$.}

\cauTN[id=T10-C1-B1-D01-VDC-092,muc=VDC,loigiai={Xét số dư của $n$ khi chia cho $7$: $P(n)$ đúng khi $n$ chia $7$ có dư thuộc tập $\{2;4\}$. Đếm trong các số từ $1$ đến $66$ được $19$ số.}]{Cho mệnh đề chứa biến $P(n)$: ``$n^2+n+1$ chia hết cho $7$''. Số các số nguyên dương $n\le 66$ để $P(n)$ là mệnh đề đúng là}
  {$19$.}
  {$20$.}{$18$.}{$22$.}

\cauTN[id=T10-C1-B1-D01-VDC-093,muc=VDC,loigiai={Xét số dư của $n$ khi chia cho $5$: $P(n)$ đúng khi $n$ chia $5$ có dư thuộc tập $\{4\}$. Đếm trong các số từ $1$ đến $89$ được $18$ số.}]{Cho mệnh đề chứa biến $P(n)$: ``$n^2+7n+1$ chia hết cho $5$''. Số các số nguyên dương $n\le 89$ để $P(n)$ là mệnh đề đúng là}
  {$18$.}
  {$19$.}{$17$.}{$21$.}

\cauTN[id=T10-C1-B1-D01-VDC-094,muc=VDC,loigiai={Xét số dư của $n$ khi chia cho $9$: $P(n)$ đúng khi $n$ chia $9$ có dư thuộc tập $\{1;4;7\}$. Đếm trong các số từ $1$ đến $94$ được $32$ số.}]{Cho mệnh đề chứa biến $P(n)$: ``$n^2+7n+1$ chia hết cho $9$''. Số các số nguyên dương $n\le 94$ để $P(n)$ là mệnh đề đúng là}
  {$32$.}
  {$33$.}{$31$.}{$35$.}

\cauTN[id=T10-C1-B1-D01-VDC-095,muc=VDC,loigiai={Xét số dư của $n$ khi chia cho $5$: $P(n)$ đúng khi $n$ chia $5$ có dư thuộc tập $\{1\}$. Đếm trong các số từ $1$ đến $72$ được $15$ số.}]{Cho mệnh đề chứa biến $P(n)$: ``$n^2+8n+1$ chia hết cho $5$''. Số các số nguyên dương $n\le 72$ để $P(n)$ là mệnh đề đúng là}
  {$15$.}
  {$16$.}{$14$.}{$18$.}

\cauTN[id=T10-C1-B1-D01-VDC-096,muc=VDC,loigiai={Xét số dư của $n$ khi chia cho $9$: $P(n)$ đúng khi $n$ chia $9$ có dư thuộc tập $\{2;5;8\}$. Đếm trong các số từ $1$ đến $39$ được $13$ số.}]{Cho mệnh đề chứa biến $P(n)$: ``$n^2+2n+1$ chia hết cho $9$''. Số các số nguyên dương $n\le 39$ để $P(n)$ là mệnh đề đúng là}
  {$13$.}
  {$14$.}{$12$.}{$16$.}

\cauDS[id=T10-C1-B1-D01-VDC-097,muc=VDC,loigiai={Với $n$ chẵn, $n^2+3n+3$ có cùng tính chẵn lẻ với $3$; với $n$ lẻ, biểu thức có cùng tính chẵn lẻ với $1+3+3=7$. a) Sai. b) Sai. c) Sai. d) Đúng.}]{Trong giờ học, cô giáo của lớp 10A5 nêu mệnh đề chứa biến $P(n)$: ``$n^2+3n+3$ là số chẵn'' với $n\in\mathbb{N}$ và đề nghị bạn Hoa phân tích tính chẵn lẻ của biểu thức theo $n$. Xét tính đúng sai của các khẳng định sau:}{%
  \y{$P(1)$ là mệnh đề đúng.}{s}
  \y{Mệnh đề ``$\forall n\in\mathbb{N}$, $P(n)$'' là mệnh đề đúng.}{s}
  \y{Mệnh đề ``$\exists n\in\mathbb{N}$, $P(n)$'' là mệnh đề đúng.}{s}
  \y{Trong các số tự nhiên từ $1$ đến $20$ có đúng $0$ số $n$ để $P(n)$ đúng.}{d}}

\cauDS[id=T10-C1-B1-D01-VDC-098,muc=VDC,loigiai={Với $n$ chẵn, $n^2+4n+1$ có cùng tính chẵn lẻ với $1$; với $n$ lẻ, biểu thức có cùng tính chẵn lẻ với $1+4+1=6$. a) Sai. b) Sai. c) Đúng. d) Đúng.}]{Trong giờ học, cô giáo của lớp 10A2 nêu mệnh đề chứa biến $P(n)$: ``$n^2+4n+1$ là số chẵn'' với $n\in\mathbb{N}$ và đề nghị bạn Hương phân tích tính chẵn lẻ của biểu thức theo $n$. Xét tính đúng sai của các khẳng định sau:}{%
  \y{$P(1)$ là mệnh đề sai.}{s}
  \y{Mệnh đề ``$\forall n\in\mathbb{N}$, $P(n)$'' là mệnh đề đúng.}{s}
  \y{Mệnh đề ``$\exists n\in\mathbb{N}$, $P(n)$'' là mệnh đề đúng.}{d}
  \y{Trong các số tự nhiên từ $1$ đến $20$ có đúng $10$ số $n$ để $P(n)$ đúng.}{d}}

\cauTLN[id=T10-C1-B1-D01-VDC-099,muc=VDC,loigiai={Số bội của $4$: $13$; của $3$: $18$; của $12$: $4$. Số thẻ thắng $=13+18-2\cdot4=23$.}]{Một trò chơi trong giờ ngoại khoá: mỗi bạn rút một tấm thẻ ghi một số nguyên dương $n$ từ $1$ đến $54$ (mỗi số xuất hiện đúng một lần). Bạn rút được thẻ sẽ thắng nếu mệnh đề ``$n$ chia hết cho $4$ hoặc $n$ chia hết cho $3$, nhưng $n$ không chia hết cho cả hai số đó'' là mệnh đề đúng. Hỏi có bao nhiêu tấm thẻ giúp người rút thắng cuộc?}{23}

\cauTLN[id=T10-C1-B1-D01-VDC-100,muc=VDC,loigiai={Số bội của $3$: $23$; của $5$: $14$; của $15$: $4$. Số thẻ thắng $=23+14-2\cdot4=29$.}]{Một trò chơi trong giờ ngoại khoá: mỗi bạn rút một tấm thẻ ghi một số nguyên dương $n$ từ $1$ đến $70$ (mỗi số xuất hiện đúng một lần). Bạn rút được thẻ sẽ thắng nếu mệnh đề ``$n$ chia hết cho $3$ hoặc $n$ chia hết cho $5$, nhưng $n$ không chia hết cho cả hai số đó'' là mệnh đề đúng. Hỏi có bao nhiêu tấm thẻ giúp người rút thắng cuộc?}{29}
